\documentclass[11pt, letterpaper]{article}

% --- TYPOGRAPHY & LAYOUT ---
\usepackage[margin=1in]{geometry} 
\usepackage[T1]{fontenc}
\usepackage[utf8]{inputenc}
\usepackage{newtxtext, newtxmath} % Superior Times Roman clone (used in scientific journals)
\usepackage[protrusion=true, expansion=true, kerning=true]{microtype} % The "secret sauce" for professional spacing
\usepackage{setspace}
\singlespacing

% --- SECTION STYLING ---
\usepackage{titlesec}
% Style sections: Large, Bold, Small Caps, with tight spacing
\titleformat{\section}
  {\large\bfseries\scshape} 
  {\thesection.}{0.5em}{}
\titlespacing*{\section}{0pt}{1.5em}{0.5em}

\titleformat{\subsection}
  {\bfseries}
  {\thesubsection}{0.5em}{}
\titlespacing*{\subsection}{0pt}{1em}{0.3em}

% --- GRAPHICS, TABLES & CODE ---
\usepackage{graphicx}
\usepackage{booktabs} % Professional tables
\usepackage{array}
\usepackage[table,xcdraw]{xcolor}
\usepackage{float}
\usepackage{caption}
% Caption style: Small font, bold label, slightly separated from figure
\captionsetup{font={small, stretch=1}, labelfont=bf, labelsep=period}

% --- HYPERLINKS ---
\usepackage{hyperref}
\hypersetup{
    colorlinks=true,
    linkcolor=black, % Keep internal links black for cleanliness
    citecolor=midnightblue, % Subtle color for citations
    urlcolor=midnightblue,
    pdfborder={0 0 0}
}
\definecolor{midnightblue}{RGB}{25, 25, 112}

% --- BIBLIOGRAPHY ---
\usepackage[backend=biber, style=authoryear]{biblatex} % Or use 'style=apalike', etc.
\addbibresource{references.bib}

% --- HEADER/FOOTER ---
\usepackage{fancyhdr}
\pagestyle{fancy}
\fancyhf{} 
\renewcommand{\headrulewidth}{0pt} % Remove header line
\cfoot{\thepage} % Center page number at bottom

% ==========================================
% DOCUMENT START
% ==========================================
\begin{document}
\newcommand{\origin}{(x_1^0,x_2^0)}
\newcommand{\landing}{(x_1^l,x_2^l)}
\newcommand{\coord}{(x_1, x_2)}
\newcommand{\shotdec}{\mathbf{x}}
\newcommand{\aimpoint}{\theta}
\newcommand{\out}{\mathbf{x}}
\newcommand{\Out}{\mathbf{X}}
\newcommand{\Samp}{\mathbf{Z}}
\newcommand{\samp}{\mathbf{z}}
\newcommand{\gridexplicit}[j]{(g^j_1, g^j_2)}
\newcommand{\grid}{\mathbf{g}}




% -------------------------------------------------------
% BLINDED TITLE PAGE
% -------------------------------------------------------
\begin{titlepage}
    \centering
    \vspace*{2.5cm}
    
    % Elegant Horizontal Rule
    \rule{\textwidth}{1.5pt}\vspace{0.5cm}
    
    {\huge \bfseries \scshape On Par Golf} \\
    \vspace{0.3cm}
    {\Large \itshape Leveraging Gaussian Process Regression for Perfect Golf} % Optional Subtitle
    
    \vspace{0.5cm}
    \rule{\textwidth}{1.5pt}
    
    \vspace{3cm}
    
    \section*{Abstract}
    \centering
    \noindent
       Golf is fundamentally a problem of sequential decision-making under uncertainty, where a surplus of strategic options often obscure the optimal path to the hole. This project presents a data-driven framework that quantifies player-specific optimal strategies by leveraging Gaussian Process Regression (GPR). Moving beyond traditional biomechanical analysis, we model the golf course as a continuous risk surface, learning the ``scoring potential'' of every location based on course geometry, historical data, and the player's shot tendencies.

       The primary output of this framework is a tangible, actionable strategy: a specific optimal club and aimpoint recommendation for any coordinate on a given hole. By simulating thousands of shot outcomes against our risk surfaces, the model solves for the decision rule that best satisfies the player's objective. Whether the goal is minimising the expected strokes to hole out or maximising the probability of securing a birdie, our framework is capable of providing a mathematically grounded roadmap, replacing intuition with calculated risk management in order to optimise performance.
    \vfill
    \today
\end{titlepage}

% -------------------------------------------------------
% MAIN BODY
% -------------------------------------------------------
\newpage
\setcounter{page}{1} % Reset count so Introduction is Page 1

\section{Introduction and Research Question / Background \& significance}

Golf is a deceptively simple sport: get the ball into the hole in the fewest number of shots possible. Upon closer inspection, however, this objective is marred by complexity. Over the course of a round, a player faces 18 holes made up of differing lengths and topographies, riddled with obstacles: water hazards, sand traps, and out-of-bounds limits with associated penalties. Consequently, no player will ever face the exact same shot twice.

Traditionally, golf performance has been viewed through a biomechanical lens. Striking a ball smaller than 2 inches in diameter with a club head moving at 135 mph toward a distant target is no small feat. Consequently, the status quo is an obsession with technical perfection: the pursuit of optimising the biomechanical motion of the swing to maximise energy transfer and consistency. This fixation is evidenced by the proliferation of AI-driven swing analysis tools such as Sportsbox.AI, GolfFix, and Sparrow Golf.

Said focus on mechanics is not statistically unfounded. In 2014, Mark Broadie provided statistical evidence that driving distance is a key factor and significant predictor of success on the PGA Tour - often more important than accuracy or putting (which had long been hypothesised to be the most important part of the golf game) \parencite{Broadie_2014a}. PGA Tour professionals who drive their ball 20 yards further than their peers typically save about 0.75 strokes per round, translating to a significant 3 strokes over a four-day tournament. The only way to hit the ball further is to maximise “smash factor”, which implies optimising technique. 

While mechanical questions are valid, they give rise to critical counter-questions: is there a way to optimise the tools we already have? How can we improve our scores without necessarily “getting better” at swinging a club? 

Ideally, the tour professional's swing is mechanically optimised, yet scoring divergences persist between the top 50 players, consistent winners, and the rest. For the recreational amateur, while technique has a large margin for improvement, meaningful mechanical changes require time, patience, and financial investment. Both professionals and amateurs can benefit from taking a different approach.

Golf is not solely about physical execution; it is a problem of complex sequential decision-making under uncertainty. There are a million sequences of different shots that can get the ball into the bottom of the hole. A player faces the course with up to 14 clubs and the possibility of hitting dozens of shot types. Such flexibility creates a ``paradox of choice''; a psychological hurdle where an abundance of options obscures the best path forward \parencite{schwartz2015paradox} and where decision theory and statistics become relevant. By treating every shot as a mathematical calculation of risk versus reward, we can mathematically determine the optimal path to the hole. 

To this end, this paper proposes a data-driven framework allowing an individual player to quantify their optimal strategy on a given golf course by leveraging both venue- and player-specific data. We 
approach this decision theory scenario through a Bayesian lens, using prior performance data to help inform future strategic choices. By implementing Gaussian Process Regression (GPR), we construct continuous predictive surfaces that evaluate the "scoring potential" of all possible landing zones on a golf course. This model allows us to mathematically define the player's objective through different "loss functions" - whether minimising the expected number of strokes to hole out for consistent play, or maximising the raw probability of a birdie when the need to exhibit a more aggressive play style arises. Ultimately, our framework replaces guesswork with calculation, identifying the specific club and aimpoint combination that statistically optimises a player's outcome from any location on the course.

The remainder of this paper is organised as follows: Section \ref{sec:Background} places this work within the broader context of sports analytics and reviews existing literature. Section \ref{sec:Data} describes the data we used, with Section \ref{sec:Modelling} exploring the modeling mechanisms employed. Section \ref{sec:Results} presents the results of our simulations, and Section \ref{sec:Discussion} discusses the limitations of the current framework alongside potential avenues for future expansion.  Appendix \ref{sec:Appendix} presents the data generating mechanism we used to supplement our existing data and Appendix \ref{sec:Appendix2} includes a glossary of relevant terms.

\section{Background and Significance}
\label{sec:Background}

The shift toward objective strategy in sports is not unprecedented. This project draws direct inspiration from SABRmetrics: the empirical analysis of baseball that uses statistics to evaluate player performance and inform objective, data-driven decisions. Coined by Bill James, this non-traditional statistical investigation famously allowed the Oakland Athletics to remain competitive against teams with significantly larger budgets by using data to identify and acquire undervalued players \parencite{moneyball_lewis}. Here, we apply the same philosophy to golf: valuing the statistical probability of a shot over the traditional ``eye test.''

Historically, golf analysis relied heavily on total score and basic counting stats (fairways hit, putts per round), which often failed to capture the true quality of a shot. This landscape changed with the introduction of Strokes Gained (SG) by Mark Broadie \parencite{Broadie_2014a}. The SG metric redefined performance evaluation by establishing baselines for the expected number of shots to hole out from any distance and lie on the course. For the first time, players had access to interpretable metrics that could isolate specific strengths and weaknesses in their game relative to the field.

Strokes Gained paved the way for applied course strategy systems, most notably DECADE \parencite{decade_golf}. DECADE is a proprietary course-management system built largely on generalised SG estimates, allowing players to make informed decisions grounded in statistical insight rather than emotion. The efficacy of such systems is well-documented; for instance, Will Zalatoris rapidly ascended from a World Amateur Golf Ranking (WAGR) of 3,000+ to the top 3 in under three months after adopting the DECADE pedagogy \parencite{decade_golf}.

The application of statistics to golf strategy is still an evolving field. Previous research has primarily focussed on three distinct methodologies: Markov Decision Processes (MDPs), neural network classification and reinforcement learning. While these studies provide a solid foundation for objective decision making, they each exhibit limitations that this project seeks to address.

Stauffer and Guillot present a study on strategy optimisation in Stroke Play for professional golfers. Their main goal was to demonstrate that a data-driven MDPs could be solved exactly for large-scale, real-world instances without relying on heuristics or Q-learning approximations \parencite{stauffer_golf_2025}. That is, they could solve the golf strategy problem perfectly for every point on the course without implementing approximations.

They inferred personalised ``TrackMan profiles'' for PGA Tour players using historical ShotLink data and constructed a 2D simulation model of a course in \texttt{C++}. By defining states as pixels and actions as triplets of (state, targeted distance, direction), they applied a value iteration algorithm to solve the stochastic shortest path (SSP) problem exactly. The model successfully demonstrated computational feasibility, allowing for the automated construction of strategic booklets outlining optimal shots from any position.

The study was limited by significant simplifying assumptions: it treated greens as homogeneous flat surfaces, assumed that in historical data players were always aiming at the pin, and always assumed risk-neutral decision-making, often leading to simulated strategies that were riskier than those observed in professional play.

Khazaeli and Javadpour focussed on implementing a probabilistic classification model to predict approach shot outcomes incorporating environmental factors using multi-level neural networks \parencite{khazaeli_golf_2024}. They combined over 700 shots from NCAA Division I golfers, and augmented them via Monte Carlo simulation. They then trained a model on 27 features, including distance, terrain, and atmospheric conditions such as wind strength and lie. They validated the hypothesis that incorporating environmental features leads to higher predictive accuracy than solely relying on distance (which still remained the dominant attribute). 

Despite its predictive power, this study was limited in scope, focussing solely on approach shots and relying heavily on synthetic data augmentation. Furthermore, it functioned primarily as a classification tool used to predict outcome zones rather than providing overall strategy optimisation.

Sugawara, Kawamura, and Suzuki explored skill level and strategy through reinforcement learning. They applied Q-learning within MDP frameworks to minimise the expected score for a hole.

They defined golfer skill via shot variance and utilised conditional probability distributions to account for penalties from trees and hazards. By comparing strategies for skilled versus unskilled agents on a simulated Augusta National, they provided the key insight that optimal strategy is highly dependent on skill level. Their simulation, however, lacked realism, the putting model was oversimplified, ignoring the slopes of the green, calculating the expected number of putts based on distance to the hole. This failure to model crucial topographical meant that their strategies were optimised for two-dimensional versions of golf, rather than the three-dimensional reality of the game.

Current literature largely relies on generalised estimates, rigid assumptions regarding course geometry, or limited subsets of the game (e.g., only modelling approach). This project addresses these gaps through a novel methodological framework. As an individual sport, golf requires player-specific solutions. We move past generalised estimates in commercial systems by integrating course-specific data and individual player tendencies into our model. Unlike previous studies that generally optimised for a single metric (lowest score), this framework implements variable loss functions. This allows for strategy optimisation based on different definitions of success, such as minimising expected strokes (conservative) versus maximising birdie probability (aggressive). Uniquely, our project leverages Gaussian Process Regression (GPR). Previously unexplored in this domain, we shift the analysis to the Bayesian paradigm, allowing us to mathematically compute the confidence intervals of our strategies, and offering a measure of certainty regarding suggested strategies that frequentist approaches lack.

\section{Data}
\label{sec:Data}

In order to understand the data inputs that we used when modelling the optimal golf strategy, it is valuable to distinguish between the ideal dataset (which would drive a fully deployed version of our tool) and the proxy dataset which we actually used in this proof-of-concept. Our model is agnostic to its data source; it requires three specific layers of information to work: course geometry for lie classification, location-specific outcome data, and player shot distribution data.

\subsection{The Ideal Data} 
\label{ref:idealData}
In a fully realized application, such as one using the PGA Tour's ShotLink system and a user's shot history, our model would consume the following:
\begin{enumerate}
    \item \textbf{Course Geometry for Lie Classification:} A high resolution, bird's-eye view vector map of the hole (where the tee is at $(0,0)$). This allows our model to embed course geometry into the model by instantly classifying any coordinate $(x_1,x_2)$ as a specific lie type (i.e. Fairway, Rough, Deep Rough or Bunker), triggering the appropriate modelling and decision logic.
    \item \textbf{Location-Specific Outcome Data:} A large repository of historical shots for a hole with a given position. For any recorded shot from a given starting coordinate, $\origin$ which can be thought of as the origin in the Cartesian plane, we would observe the discrete outcomes of how many shots it took players to hole out to the corresponding hole. The numbers of shots would encode the underlying ``difficulty'' of a specific spot: locations with the same Euclidean distance to the pin where the average strokes to hole out are higher than others are inherently ``harder''. \textcolor{blue}{insert drawing}
    \item \textbf{Player Shot Distribution Data} A library of ``stock shots'' labeled by the player. Whether recorded on a launch monitor (e.g., Trackman, GCQuad) or manually logged, the dataset would consist of the landing coordinates $\landing$ generated by hitting shots from  $\origin$ aiming down a set target line (where the lateral offset is zero) for every repeatable shot in a player's repertoire\footnote{In addition to the 14 clubs they usually carry in a bag, competitors usually identify different kinds of shots they can hit with their clubs. These shot-types tend to impact the ball's flight affecting factors such as the ball's lateral curve, apex height, or total carry distance. While most players recognise and play standard full swing shots, a player may choose to name different shots in distinct ways (e.g., a player may name a specific shot an L-shot, while another may name it a 3/4 swing). Trackman allows the player to record and save shots under different names, which are then modelled in such a way that the output of our framework includes recommendations implementing such labelled shots.}. 
    
    Ideally, we would have labelled data from different starting lies such as rough or fairway bunkers, accounting for the increased dispersion and variance that results of hitting from different surfaces. More concretely, the player's shot library would contain observations encoding the tuple $\left(\origin, \;\landing,\; Lie\right)$ for every specific labelled $Shot$ type. From data in said format, we derive enough information to fully characterise the conditional landing distribution $P\left(\landing\;|\;Shot,\; \origin,\; Lie\right)$.
\end{enumerate}

\subsection{Data Implemented in this Study}

To demonstrate the viability of our framework, we developed and implemented a robust proxy dataset that provided an appropriate substitute for the three core data layers required by the model \textcolor{blue}{(should i again refer the reader to the Appendix?)}. We successfully digitised a real golf course using satellite imagery and vector mapping tools, achieving the fidelity of the ideal data for our course geometry. In absence of proprietary PGA Tour ShotLink data, we built a data-generating mechanism to emulate location-specific strokes-to-hole-out outcomes, supplemented by empirical estimates. Finally, we established a representative LPGA\footnote{LPGA} Tour player profile using publicly available launch monitor statistics to model the player's shot distribution data.

\subsubsection{Course Geometry}
We digitised hole 9 of Mountain Meadows Golf Course in California, a Par 3 chosen for its rich interaction of obstacles such as bunkers and water hazards near the green. The course layout was scraped using OpenStreetMap \parencite{openstreetmap}, and refined in QGIS \parencite{qgis} to create accurate vector polygons compatible with the \texttt{Shapely} \parencite{shapely} library in Python. To test chained-decision-making (Par 4 strategy), we manually generated a hypothetical extension of hole 9 adding a secondary water hazard to increase strategic complexity. \textcolor{red}{insert images below}

\subsubsection{Emulating Location-Specific Outcome Data}
\label{subsec: greendata}
\begin{itemize}
    \item \textbf{On the green:} We generated a dense sample of synthetic putting outcomes for our digitised green taking into consideration slope and distance in a three dimensional setting. The data generated simulates the ``ideal'' data described in subsection \ref{ref:idealData}, providing a count of the strokes to hole out for hundreds of points on the putting surface. \textcolor{red}{image}

    \item \textbf{Off/Around the Green:} Generating realistic outcomes around the green, while having the potential for being extremely informative for our model, was deemed infeasible (due to the lack of available estimates, and the variance of outcomes and approaches that different players may take when tackling short range chips or bunker shots). We therefore interpolated average ``strokes to hole out'' as a function of distance values from different lies from Mark Broadie's empirical estimates (derived from the PGA Tour Shot Link data from 2003-2012) and published in  Every Shot Counts \parencite{Broadie_2014a}. These estimates can also be found in the Appendix \ref{app: short-game data} and allow us to assign a baseline expected value to any lie around the green without needing a generated point cloud for that specific coordinate.
    
\item \textbf{Long Game:} We opted not to generate granular outcome distributions for longer shots due to the high computational cost. The optimality of any long shot is fundamentally determined by the quality of the resulting landing state (i.e., the expected strokes to hole out from the new position. The primary objective of the long game is to transition to a desirable short-game state, making the precise outcome distribution of the shot itself secondary to the value of the resulting location. Explicit shot-by-shot was deemed unnecessary, and as such, chose to rely on the modelling logic and priors of our results to approximate the distributions for these shots. We outline this process in Section \ref{subsec: strategicBlanket}.
\end{itemize}

\subsubsection{Player Shot Distribution Data}

We utilised a proxy shot distribution data set based on publicly available Trackman launch monitor statistics for a representative ``average'' LPGA tour player. The library encompassed 13 standard clubs from $60^\circ$ wedge to a Driver (ranging approximately $66$ to $250$ yards), plus several partial wedges for shorter distances. For each of the $20$ shot-types recorded in this proxy player's repertoire, the dataset contains 30 recorded landing coordinates, totalling 600 observations. We assume all of these shots are hit off the fairway (or a tee in the case of a Driver). These coordinates were generated from shots hit along a single target line (zero lateral offset), providing the necessary dispersion characteristics to model the player's skill profile. 

\textcolor{red}{image}

\section{Modelling and Methods}
\label{sec:Modelling}

\subsection{Our Mathematical Framework: Decision Theory}

Ben Hogan, who was widely regarded as one of the greatest ball strikers in golf history, famously stated that even on a great day, he would hit at most three  truly great shots in a round \parencite{Jenkins_1998}. The inherent imperfection in performance means golf fundamentally poses a problem of decision-making under uncertainty.

In statistical decision theory, we usually seek a decision rule $\delta$ that minimises a specific loss function $\mathcal{L}$. In the context of a golf shot, our decision rule is a tuple $\delta = (s, \aimpoint)$, where $s$ represents the discrete choice of club/shot type and $\aimpoint \in \mathbb{R}$ represents the continuous choice of aimpoint's lateral offset from the pin (where negative values denote aiming left of the pin, and positive values denote aiming right). The outcome of this decision, which is the actual landing coordinate of the ball, is denoted by the random vector $\Out \in \mathbb{R}^2$. Given that human performance is imperfect and that there are ever-changing variables, the landing is stochastic, drawn from a probability density function conditioned on the player's decision: $\Out \sim f(\out | s, \aimpoint) $. We denote the realised outcome of this variable as $\out^l = \landing$.

The objective of our model is to find the decision $\delta$ minimising the expected cost of a shot. The loss function, $\mathcal{L}(\Out)$, measures the cost (or penalty, in strokes) incurred when the ball lands at coordinate $\Out$. In a deterministic scenario where the outcome $\Out$ was certain, we would simply minimise the known loss $\mathcal{L}(\Out)$. However, we are operating in an imperfect and therefore stochastic paradigm, where the outcome is governed by the probability density function $f(\out | \delta)$. The nature of golf requires minimising the risk, $R(\delta)$, defined as the expected loss over the player's shot outcome distribution:

$$R(\delta) = \mathbb{E}_\Out[\mathcal{L}(\Out)] = \int_{\mathbb{R}^2}\mathcal{L}(\out)f(\out | \delta)d\out$$

The choice of loss function, $\mathcal{L}$, critically shapes the optimal strategy. In standard statistical modeling, we often select symmetric loss functions. For example, the Squared Error Loss, $\mathcal{L}^*(P, \hat{P}) = (P - \hat{P})^2 $, used to derive the mean, assumes that an error of magnitude $\epsilon$ is equally detrimental regardless of direction. In these general cases, the loss is a function of the true parameter being estimated, $P$, and the chosen predictor, $\hat{P}$.

However, golf, like many real-world scenarios often presents asymmetric loss landscapes. Consider civil engineers designing bridges: building a bridge a foot too long is a minor inefficiency and cost when compared to the catastrophic failure of building it a foot too short. Golf exhibits similar non-linearities. A shot missing a couple yards to the left might rest safely in a spot fractionally less optimal than the ``perfect'' spot, whereas one missing a couple yards right could plunge into a water hazard, incurring a fixed penalty stroke and a positional disadvantage.

While we can conceptually define these loss scenarios, the resulting risk function $R(\delta)$ remains an elusive unknown surface over the two dimensional domain of a golf course. We don't have an analytical formula that outputs the expected loss for any given coordinate $\out^l$: we only have sparse, noisy observations of shots and scores. To perform any sort of optimisation, we need to first learn the associated risk surface for a specific hole (no two holes are exactly alike!). The definition of such a surface requires a statistical method capable of taking in scattered data to predict a continuous landscape of expected strokes to hole out (ESHO), a task for which we can leverage Gaussian Process Regression (GPR). 

In this project, we define two distinct risk landscapes that our GPR models can learn, corresponding to different player objectives.
\begin{enumerate}
    \item \textit{Minimising Expected Strokes to Hole Out.} Here, our loss function, $\mathcal{L}_\mu$,  is the count of strokes to hole out from any given landing position. $ \out \in \mathbb{R}^2$, and outputs an integer count of strokes.

    The function is defined as: $\mathcal{L}_{\mu}:\; \mathbb{R}^2 \to \mathbb{Z^+}$
    $$\mathcal{L}_\mu(\out^l)= \text{Strokes to hole out from }\out^l$$
    The associated risk, $R_{\mu}$, is the expectation if the aforementioned loss given a specific decision $\delta$. The risk function domain is consistent with the decision rule:  $R_{\mu}:\; \mathcal{S} \times \mathbb{R}\to \mathbb{R}$, where, $\mathcal{S}$ is the discrete set of shot types a player chooses from.

    $$R_{\mu}(\delta) = \mathbb{E}[\mathcal{L}_\mu(\Out) | \delta]$$
    
    By minimising $R_{\mu} (\delta)$ we identify decisions where the average outcome (Expected Strokes to Hole Out) is the lowest. We naturally account for asymmetric hazards because the penalised observations fundamentally skew the local mean upwards. The actual optimisation mechanism is explored in Section \ref{subsec:par3} and \ref{subsec: par4}.

    \item \textit{Maximising Birdie Probability.} 
    
    Here, we employ a utility-based loss function, where ``success'' and, by complementation, loss, is binary. The loss is defined based on whether the expected score from the landing spot $\out$ fails to achieve the birdie target. 

    The loss function $\mathcal{L}_B$, maps the landing position $\out$ to a binary outcome, $\mathcal{L}_B: \mathbb{R}^2 \to \{0,1\}$.

    $$\mathcal{L}_B(\mathbf{x}) = \begin{cases} 0 & \text{if } \text{score from } \mathbf{x} < \text{Par Score, (i.e. Birdie or better)} \\ 1 & \text{if } \text{score from } \mathbf{x} \ge \text{Par Score (i.e. Par or worse)} \end{cases}$$

    The associated risk, $R_B$, is the probability of failure (not making a birdie or better), $R_B: (\mathcal{S} \times \mathbb{R}) \to [0,1]$. Since the loss is binary, the expectation is equivalent to the probability of the loss occurring: 
    
    $$R_B(\delta) = \mathbb{E}[\mathcal{L}_B(\Out)|\delta] = \mathbb{P}\left(\text{Strokes to Hole Out from } \Out \geq \text{Par Score} | \delta \right)$$
    
    The minimisation of this value reflects the ``must-score'' scenarios where the magnitude of a miss is irrelevant (a par is as useless as a triple bogey). 

     The actual optimisation mechanism is explored in Section \ref{subsec: par4b}.
\end{enumerate}

\subsection{Our Modelling Engine: Gaussian Process Regression}

\subsubsection{Gaussian Process Regression to Learn Risk in Golf}
In order to minimise the risk functions defined above, helping us find the optimal decision rules, we have to first estimate their outcomes. We treat the unknown risk surfaces (expected strokes or birdie probability) corresponding to different holes as latent functions $f(\out)$ to be learned from the observed data. An ideal mechanism to estimate the risk surface is Gaussian Process Regression.

Gaussian Process Regression is a tool that works intuitively with this project because golf courses are inherently spatial environments where nearby locations (on the same kind of surface) tend to have similar strategic values. If a player is 100 yards from the hole in the fairway, moving a couple yards on the same surface is unlikely to drastically change their expected score. GPR's smoothness asumptions align with the green, allowing us to interpolate values across unobserved regions based on spatial proximity.

Furthermore, golf scoring relationships are complex and non linear. A putt from 10 feet is not necessarily twice as hard as one from 5 feet, nor do two different 100-yard shots share the same difficulty if one is from deep rough and the other is off the fairway. Unlike rigid parametric models (such as linear regression) which might force specific functional forms, GPR is non-parametric. It lets the data speak for itself, ``learning'' arbitrary functional relationships and adapting to the true complexity of the three dimensional course surface.

Lastly, GPR distinguishes between high-confidence predictions (where data is dense) and uncertain regions (where data is sparse or highly volatile) through its posterior variance. When asking questions about risk management, knowing our level of confidence in a strategy may be just as important as the strategy itself.

\subsubsection{What are Gaussian Processes and what is Gaussian Process Regression?}

Gaussian Processes are collections of random variables where any finite subset follows a multivariate Gaussian distribution. As expected, they are fully specified by their mean function $m(\out) = \mathbb{E}(f(\out))$  and covariance kernel $k(\out,\out') = \mathbb{E}[(f(\out)-m(\out))(f(\out')-m(\out'))]$:

$$f(\out) \sim \mathcal{GP}(m(\out), k(\out,\out'))$$

In our context, $\out$ would represent the spatial coordinates on the hole ($\out = \coord$), and $f(\out)$ would represent the latent scoring potential of any given position (i.e., f encodes the number of shots to hole out from a given position).

\subsubsection{Covariance Kernels}
The heart of Gaussian Processes are the kernel functions, $k(\out,\out')$ which encode our prior assumptions about the structure of the data. To capture the aforementioned spatial coherence of a golf course, we use the Radial Basis Function (RBF) kernel:

$$k(\out,\out') = \sigma^2 \exp\left(\sigma^2-\frac{||\out-\out'||^2}{2l^2}\right)$$

Where the hyperparameters dictate the surface's behaviour and flexibility. The length-scale, $l$, controls the flexibility of our functions controlling smoothness by dictating how quickly scoring difficulty changes over distance. Smaller values of $l$ drive correlations to 0, allowing functions to exhibit ``wilder'' or more jagged behaviour, while also relaxing the width of confidence bounds. The signal variance, $\sigma^2$, on the other hand controls the vertical amplitude of the function we are trying to estimate.

\subsubsection{The Posterior: From Prior to Prediction}

We model our observed data $\mathbf{y}$ (i.e., actual strokes taken from a spot) as the latent function with Gaussian noise: $\mathbf{y} = f(\Out) + \epsilon$, where $\epsilon \sim \mathcal{N}(0, \sigma^2_n)$. 

Given a set of $n$ observed training points $\Out$ and outcomes $\mathbf{y}$, and a new test location $\out_*$ (a new spot we want to evaluate), the joint distribution of the observed data and the unobserved latent function value $f_*$ is given by:

\begin{align*}
    \begin{bmatrix}
        \mathbf{y}\\
        f_*
    \end{bmatrix} \sim 
    \mathcal{N}\left (\begin{bmatrix}
        k(\Out,\Out) + \sigma^2_nI & k(\Out,\out_*) \\
        k(\out_*, \Out) & k(\out_*, \out_*)
    \end{bmatrix}\right)
\end{align*}

Where $k(\Out,\Out)$ is the covariance matrix of the training points. By conditioning on the observed $\mathbf{y}$, we derive the posterior predictive distribution for the risk at our new location $\out_*$:

\[f_* |\Out, \mathbf{y}, \out_* \sim \mathcal{N}(\bar f_*, \text{Var}(f_* | \Out, \mathbf{y}, \out_*)\]

Where the posterior mean $\bar f_*$ provides our point estimate for optimal strategy described in Section \ref{subsec:par3} (predicted ESHO or Birdie Probability) and the posterior variance quantifies our uncertainty.

GPR allows us not only to predict the mean outcome, but to understand the confidence of said predictions (essential for robust decision-making in areas suffering from sparse data density).

\subsection{Modelling an Optimal Shot. Par 3s \& Minimisation of Strokes.}

\label{subsec:par3}

With our modelling mechanism in place, we can now simulate the decision-making process for a specific shot. We begin with a Par 3, the ``simplest'' type of hole in golf, where we aim to reach the putting surface in a single stroke, and therefore don't need to chain strategic decisions. 

We model the problem as a decision taken from a fixed starting coordinate $\origin$ which we treat as the origin $(0,0)$ in the shot's local coordinate system. Our optimisation follows a `"green-to-tee" philosophy: by understanding the value of every landing spot on the green, we can inform the optimal decision from the tee (the choice that precedes it). Such an implementation is viable because the problem exhibits optimal substructure, where the ultimate solution to a problem depends on the optimal solutions to its smaller subproblems. Said characteristic is best known as Bellman's Principle of Optimality\parencite{rosu_bellman}.

\subsubsection{Step 1: Fitting GPR to the Green}
Our strategy builds on the GPR model fitted to the green data (as described in the Section \ref{subsec:par3}), evaluated at the posterior mean. The fitted Gaussian defines a continuous Expected Strokes to Hole Out (ESHO) surface $f_G: \Omega_{green} \subset \mathbb{R}^2 \to \mathbb{R}$, which outputs the expected number of strokes to hole out (ESHO) from any coordinate on the putting surface. In other words, $f_G$ directly implements our loss function $\mathcal{L}_\mu$ for the putting region.


Because $f_g$ is fit to the unique sample of data specific to hole 9's green and particular pin position, $f_G$ intrinsically accounts for the non-linear difficulties (slope, tiering) that a simple distance-based function would miss.

\textcolor{blue}{insert green gpr image}

When a simulated shot from the fairway or tee lands on the green, we compute its ``cost'' by evaluating its scoring potential from its supposed resting position.

If a shot lands on the green, the cost is retrieved directly from $f_G(\out)$. 

Alternatively, if it misses the green, we employ an interpolation function $f_A: \mathbb{R}^+ \times \mathcal{C} \to \mathbb{R}^+$, where $\mathcal{C} = \{F, R, S\}$ represents the set of lie conditions (fairway, rough or sand trap). $f_A$ is a function of lie and distance from the hole that outputs the expected number of shots to hole out. The function is derived from Broadie's empirical estimates for each lie (See Appendix \ref{app: short-game data}).

Finally, if a shot lands in a water hazard, we simulate a ``drop'' at the nearest point of relief (drawing a direct line from the balls starting point, and simulating that the ball landed at the first intersection with land), incurring a penalty stroke and recursively evaluating the cost from the new position.

Formally, we define a global, around the green cost function $AG(\out)$, representing the expected strokes to hole out from any location $\out = \coord$ near or on the green:
\[
AG(\out) = \begin{cases}
    f_G(\out^l)& \text{if } \out^l \in \Omega_{green} \\
    f_A\left(\out^l, \text{Lie}(\out^l)\right) & \text{if } \out^l \in \Omega_{fairway}\; \cup \Omega_{rough} \; \cup \Omega_{sand}  \\
    1 + AG(\out^{drop}) & \text{if } \out^l \in \Omega_{water}
\end{cases}
\]
\subsubsection{Step 2: Simulating Player Outcomes}
\label{subsec: playerSim}
In order to evaluate the success of a strategy in approaching the green, we must be able to model and simulate the stochastic nature of a player's ball striking. We achieve this by sampling from a normal distribution fitted to the player shot distribution data described in \ref{sec:idealData}, and then transforming the outcomes based on the chosen aimpoint.

\begin{enumerate}
    \item \textit{Fitting Bivariate Gaussians.} First, we analyse the player's historical shot data to create statistical profiles for each of the available shots $s$. Using standard statistical libraries such as \parencite{numpy} Numpy in Python, we empirically estimate the two-dimensional centre of the data $\mathbf{\hat\mu_s}$ and the covariance matrix $\mathbf{\hat\Sigma_s}$ for the landing coordinates of shots aimed directly at the pin (zero lateral offset) hit with shot $s$, and in this case assumed to be hit off the fairway $F$. 
     \[\Out \sim \mathcal{N}(\mathbf{\hat\mu_s}, \mathbf{\hat\Sigma_s})\]   
    This step ``learns'' a player's unique dispersion shape, capturing how their distance control and directional accuracy may vary with different shots.

     \item \textit{Monte Carlo Simulation.} We then employ a Monte Carlo simulation to generate a synthetic dataset of potential outcomes. For a selected shot $s$, we draw $N$ random samples from the fitted multivariate normal distribution. 
     $$\Samp_s = \{\samp^{s1}, \samp^{s2}, \ldots, \samp^{sN}\} \text{ where } \samp^{si} \sim \mathcal{N}(\hat\mu_s, \hat\Sigma_s)$$
     This creates a sample of shot outcomes statistically mirroring the player's tendencies.

     \item \textit{Changing Aimpoints.} Finally, in order to evaluate different strategic aimpoints, we apply a linear transformation to our simulated shots. If a player aims at an angle $\aimpoint$ relative to the pin, we rotate the entire set of simulated outcomes $\Samp_s$ using a standard rotation matrix $R(\aimpoint)$: 
     \begin{align*}
         R(\aimpoint) &= \begin{bmatrix}
             \cos \aimpoint & - \sin \aimpoint \\
             \sin \aimpoint & \cos \aimpoint
         \end{bmatrix} \\
         \samp^{si\aimpoint} &= R(\aimpoint)\samp^{si}
     \end{align*}
     Giving us $\Samp_{s\aimpoint} =  \{\samp^{s1\aimpoint}, \samp^{s2\aimpoint}, \ldots, \samp^{sN\aimpoint}\}$ and allowing us to evaluate the player's dispersion patterns at different aimpoints from $-20 ^\circ$ to $+20 ^\circ$ relative to the pin, ultimately evaluating the expected cost of rotated samples across the course's risk surface.
\end{enumerate}
\textcolor{blue}{insert sample evaluated on green image}

\subsubsection{Step 3: Minimising Expected Strokes}
\label{subsec: minimise}

We proceed to the optimisation step, which returns the optimal strategy. For a given strategy, defined by a labeled shot type $s$ and aimpoint $\aimpoint$, we have a set of $N$ simulated landing positions $\Samp_{s\aimpoint} = \{\samp^{s1\aimpoint}, \samp^{s2\aimpoint}, \ldots, \samp^{sN\aimpoint}\}$. 

For each simulated shot $i$ (which has a defined $\aimpoint$ and $s$), we evaluate the outcome by querying our  $AG(\out)$ at the landing coordinate $\samp^{si\aimpoint} =  R(\aimpoint)\samp^{si} $. The ``cost'' of the shot is the sum of the stroke taken to reach the first point plus the expected strokes required to hole out.

$$\text{Cost}^{i,s, \aimpoint} = 1 + AG(\samp^{si\aimpoint} )$$

\textcolor{red}{It's verbose but i like how direct it is.}

We then aggregate these individual outcomes to compute the empirical risk (average expected strokes) for the specific strategic configuration:
$$R_\mu(\delta) = R_\mu(s, \aimpoint) = \frac{1}{N} \sum ^N_{i=1} (1+ AG(\samp^{si\aimpoint}))$$


This computation is repeated iteratively across all feasible shots in the player's bag and a discretised range of aimpoints (e.g., $\aimpoint \in [-20^\circ, +20^\circ]$ in $1^\circ$ increments). The optimal strategy for the this shot and by extension, hole, is the pair $(s^*, \aimpoint^*)$ which minimises the risk surface:

$$\delta_{opt} = \arg \min_{s, \aimpoint} R_\mu(s, \aimpoint) = (s^*, \aimpoint^*)$$

We convert the optimal angle $\aimpoint^*$ to yards (left or right of the pin) by using basic trigonometric conversions that make the output actionable for the player. Our three-step process transforms abstract data into concrete recommendations that the player can then operate on.

\textcolor{blue}{add image of output}

\subsection{Modelling Chained Decisions. Par 4s \& Recursive Optimisation.} 
\label{subsec: par4}
Extending on the innovative use of GPR modelling on Par 3s, we apply GPR to longer holes requiring chained decision-making. Our optimisation strategy for Par 4 holes continues the recursive ``green-to-tee'' philosophy. 

The key insight is that we believe Par 4s are Par 3s with a variable teeboxes when approaching the green (that is, we need to consider different positions we may be presented with, when approaching the green). To optimise the actual tee shot, we first understand the strategic value of every possible intermediate landing zone in the fairway, rough or sand. We achieve this through a two-stage process: pre-computing the optimal approach from all possible positions and then fitting a GPR surface to those optimal values.

\subsubsection{Step 1: The "Strategic Blanket" (Pre-computing Approaches)}
\label{subsec: strategicBlanket}
We begin by evaluating a dense grid of potential second-shot locations on and near the fairway, in our case ranging from $50$ to $350$ yards from the green. We define this discrete grid of points as $\mathcal{G} = \{\grid ^1, \grid^2, \ldots, \grid^M\}$, treating each point $\grid^j = \gridexplicit{j}$ as a temporary tee box.

From each point $\grid^j$, we run the complete Par 3 simulation algorithm described in the previous in Section \ref{subsec: minimise}.

\begin{enumerate}
    \item Identify the lie at $\grid^j$ - fairway, rough, or bunker –  to select the appropriate shot distribution to sample from the player shot distribution dataset. If $\grid^j$ is water, we simulate dropping with the same logic described in the previous section.
    \item Simulate shots for all clubs and aimpoints as described in Section \ref{subsec: playerSim} from this temporary ``origin''.
    \item Identify the strategy $(s^*, \aimpoint^*)$ that minimises expected strokes from that specific location.
    \item Report the minimum risk value (ESHO) $R_{\mu}(s^*, \aimpoint^*)^j$, to use the values in the next phase.
\end{enumerate}
This process yields a minimal Expected Strokes to Hole Out (ESHO) value for every point on the grid. We then fit a \textit{second} Gaussian Process model to these ESHO values associated to each $\grid^j$. This creates a new continuous surface $f_{approach}: \mathbb{R}^2\to \mathbb{R}^+$ which we name a ``strategic blanket''. 

Evaluated at a coordinate $\out^l$ on the golf course, this function does not return a value representing the raw difficulty of the hole, but rather represents the best possible expected score a player can achieve from 50 - 350 yards from the green, assuming they play optimally from that point forward.
\textcolor{blue}{insert hole gpr image}

\subsubsection{Step 2: Optimising the Tee Shot}
With $f_{approach}(\out)$ acting as our main evaluation function (replacing $AG(\out)$), the Par 4 tee shot optimisation becomes mathematically identical to the Par 3 problem.

We simulate tee shots for long-range clubs (Driver, Woods, and Long Irons) using the player's dispersion parameters $\hat\mu_{s},\hat\Sigma_{s} $, but now assuming to be hitting off a tee. For each simulated landing point $\samp^{si\aimpoint}$ (derived from the same Monte Carlo process described in Section \ref{subsec: playerSim}), the cost is no longer derived from the green model, but from our strategic blanket:
$$\text{Cost}^{i, s, \aimpoint} = 1 + f_{approach}(\samp^{si\aimpoint})$$
Our optimisation seeks the shot type and aimpoint that lands the ball in the ``valleys'' of $f_{approach}$ – areas where the subsequent shot is easiest and therefore the expectation of the optimal choice is lowest - while avoiding the ``peaks'' (hazards or difficult angles). 

We aggregate these costs once again to find the risk of the tee shot:

$$ R _{\mu_{tee}}(\delta)= R_{\mu_{tee}}(s, \aimpoint) = \frac{1}{N} \sum^N_{i =1} (1 + f_{approach}(\samp^{si\aimpoint}))$$

As such, the optimal tee strategy is defined as:
$$\delta_{opt}  = \arg \min_{s, \aimpoint} R_{\mu_{tee}}(s, \aimpoint) =  (s^*_{tee}, \aimpoint^*_{tee})$$

This recursive framework allows the model to make nuanced trade-offs potentially suggesting more conservative options such as hitting the shorter distance club not just to ensure landing on the fairway, but specifically because the Driver may introduce other trouble (water or bunkers into play).

\textcolor{blue}{add image of final output for par 4 with good description of how it should be read}

\subsection{Modelling Alternative Loss Functions. Maximisation of Birdie Probabilities.}
\label{subsec: par4b}

Up until now, our optimisation strategy has focused on minimising the expected number of strokes to hole out. This approach naturally produces conservative strategies – decisions that prioritise avoiding high-penalty disasters (such as hitting the ball into water hazards or Out of Bounds) over chasing rare, high-reward outcomes. In statistical terms, this conservative play minimises the variance of the score distribution to protect the mean.

However, golfers are often faced with high-tension scenarios where consistency becomes inconsequential. If a player needs a birdie to make the cut, force a playoff, or win a tournament, a par may be functionally equivalent to a double bogey. In these moments, aggressive strategies are required. Here, the player willingly accepts a higher expected score (and higher variance) in exchange for increasing the probability mass in the lower tail of the distribution (i.e. in the case of a par 4, the chance of making a 3 or better).

To model this change in objective, we alter our loss function by converting the data used to train the GPR.

\subsubsection{Step 1: The Probability Surface}
We begin by returning to the raw data generated on the green (as descibed in Section \ref{subsec: greendata}. In the par 3 and par 4 expectation minimisation models described earlier, each data point collected on a green was a count of strokes (e.g., 1, 2, or 3) un a given location $\coord$. For this new objective, we perform a binary conversion on the training data outcomes $\mathbf{y}$. 
\[y_{binary} = \begin{cases}
1, & \text{if strokes} \le 1 \;\text{(holed)} \\
0, & \text{if strokes} > 1 \;\text{(missed)}
\end{cases}\]

\textit{Why do we care about "one-putts"?} A typical hole's par is set such that a typical golfer will achieve a par by taking at most two putts to hole out once on the green. Therefore, the most secure way to score a birdie is to take at most one putt to hole the ball. The likelihood of one-putting is what we are therefore trying to model here.

We then fit a new GPR model to these binary outcomes. Although the inputs are discrete, the GPR learns a continuous latent function $f_{birdie}: \mathbb{R}^2 \to [0,1] $ representing the conditional probability of one-putting from any location $out$ on the green. This creates a ``Birdie Probability Surface'', which decays exponentially: being 3 feet away from the cup offers massive reward, whereas being 15 feet away might offer near-zero value in ``birdie necessary'' circumstances.

\subsubsection{Step 2: Recursive Maximisation for Birdie Likelihood Maximisation}
With this green GPR surface defined, we re-run the original simulation mechanism described in Section \ref{subsec: minimise} with a modified evaluation logic:

\begin{enumerate}
\item \textbf{Approach Simulation:} For all simulated approach shots $\samp^{si\aimpoint}$ (the rotated samples from the Monte Carlo step), we query the new surface:

\[
\mathbb{P}(\text{Birdie} \mid \mathbf{z}^{si\aimpoint}) = \begin{cases}
    f_{birdie}(\mathbf{z}^{si\aimpoint}) & \text{if } \mathbf{z}^{si\aimpoint} \in \Omega_{green} \\
    0 & \text{else}
\end{cases}
\]

Note that we make the simplifying assumption that shots missing the green ($\samp^{si\theta} \not\in \Omega_{green}$) have negligible probability of being holed for birdie compared to putts.

\item \textbf{Risk Calculation:} The risk we seek to minimise is the probability of failure (not achieving a birdie or better). We therefore aggregate the probabilities across all $N$ samples:

$$R_{B}(\delta)= R_{B}(s, \aimpoint) = 1 - \frac{1}{N} \sum_{i=1}^{N} \mathbb{P}(\text{Birdie} \mid \mathbf{z}^{si\aimpoint})$$

\item \textbf{Optimal Aggressive Strategy on Par 3s:} Finally, we identify the strategy that minimises the probability of failure (thereby maximising birdie probability) by iterating over all feasible shots and aim angles:

$$\delta_{opt} = \arg \min_{s, \aimpoint} R_B(s, \aimpoint) = (s^*, \aimpoint^*)$$

This yields the specific club and aimpoint required maximise the odds of acquiring a birdie, accepting the higher variance associated with the shot in exchange for the maximised potential of a low score.
\end{enumerate}

\subsubsection{Step 3: Aggressive Tee Shots}

For par 4s, we repeat the ``strategic blanket'' process described in Section \ref{subsec: par4}, but with a fundamental shift in the evaluation metric.

Instead of fitting a GPR to the ESHO from every fairway location, we now map the highest birdie probability achievable from every grid point $\mathbf{g}^j$.

Alike the par 3 birdie optimisation mechanism, this creates a new surface, $f_{approach_b}: \mathbb{R}^2 \to [0,1]$, where the ``peaks'' represent fairway locations that offer the best statistical chance of converting a birdie on the subsequent shot.

When optimising the tee shot, the model aggregates these probabilities to find the strategy that maximises the potential for a birdie. We simulate the tee shots as before, but the cost function is inverted to represent the probability of success:
$$R_{B_{tee}}(s, \aimpoint) = 1 - \frac{1}{N} \sum^N_{i =1} f_{approach_b}(\mathbf{z}^{si\theta})$$

The output of this optimisation can often be drastically different from the conservative model. Which in some ways, is exactly what we want to see – different optimal decisions for different objectives. 
\section{Results}
\label{sec:Results}
The primary objective of this study was to demonstrate that the implementation of Gaussian Process Regression was feasible when trying to simulate golf performance. Ultimately we want to demonstrate that we could generate actionable strategic recommendations based on varying player objectives. Our simulation results confirm that our model effectively learns course-specific risk surfaces and adapts its recommendations accordingly.

\subsection{Risk Surface Generation}
The GPR model successfully learns continuous risk surfaces from discrete, noisy data points. For the Par 3 case study (hole 9, Mountain Meadows), the model generated a smooth Expected Strokes to Hole Out (ESHO) surface that accurately reflected the penalty structure of the hole. The surface naturally exhibited higher ``risk'' values near the water hazard and bunkers, and lower values in the wider part of the green, effectively interpolating difficulty in unobserved regions. Our model subsequently gave us a very reasonable recommendation.

\subsection{Strategic Divergence: Conservative vs. Aggressive}
The most significant finding is the clear divergence in optimal strategy when switching between loss functions (minimising expected strokes versus maximising birdie probability).

\begin{itemize}
    \item \textbf{Minimising Expected Strokes (Conservative):} When minimising $R_\mu(\delta)$, the model recommends strategies that prioritise variance reduction. On the par 4 simulation, the ``Strategic Blanket'' identified that the optimal tee shot was not necessarily the driver. Instead, our model suggested hitting a 5-iron (a very safe choice off the tee) aimed at the widest part of the fairway, sacrificing distance to avoid the high-penalty zone (the water penalty crossing the hole) that the driver's wider dispersion would bring into play. 
    
    \item \textbf{Maximising Birdie Probability (Aggressive):} Conversely, when minimising $R_B(\delta)$, the model shifted to a high-variance, high-reward approach. The model recommends a more aggressive approach by hitting the driver over the water and accepting a non-zero probability of a penalty in exchange for a significantly shorter approach shot. Which would intuitively have a smaller dispersion.
\end{itemize}

In summary, the framework works as intended: it does not simply output a single ``best'' shot, but effectively quantifies the trade-offs between consistency and scoring potential, providing players with a roadmap that they can implement regardless of the outcome of their first shot.

\textcolor{red}{Simple results section. I am thinking the caption of the images with the strategy can have all the details. I want these images in the "areas" where they were generated, so i think i should keep the description there.}

\section{Discussion}
\label{sec:Discussion}

This project serves as a proof-of-concept for the derivation of mathematically reasonable golf strategies under a Bayesian paradigm. By treating the course as a continuous probabilistic surface, we accept the variability and uncertainty behind a shot, and choose to model it. As outlined earlier, our approach works, however, several limitations must be addressed before this framework can be deployed in a professional tournament setting.

Perhaps the most significant of our limitations is the simplification of our ball flight physics. Our model currently assumes that the ball stops at its carry distance, when in reality shots usually have roll-out. How much a ball rolls out is complicated to compute as it depends on landing angle, spin rate, and ground firmness. On hard fairways or greens a ball may roll-out 20-40 yards, potentially bringing hazards into play that we deemed unreachable in our statistical analysis based on carry distance alone. Furthermore, we do not model the vertical trajectory (apex height) or shot shape (draw/fade), meaning the model cannot currently account for obstacles that require clearance, such as tall trees.

The current iteration of our project also assumes a static environment. We ignore meteorological conditions such as wind speed, direction, temperature, and air density, all of which drastically affect carry distances. While these variables can be parametrised, implementing variables that can change quickly during a tournament presents a significant challenge. Additionally, our model assumes a constant green speed (Stimp) and ignores elevation changes from tee to green, which are critical for accurate club selection.

While our green putting model is robust, the ``around-the-green'' data relies on sparse interpolated averages from historical PGA Tour baselines rather than player-specific data. This assumes the player possesses ``average'' professional scrambling ability, which may mask individual weaknesses (e.g., a player who is statistically poor from sand bunkers should be advised to avoid them more aggressively than our model currently suggests).

Our project also presents significant opportunities. A major advantage of GPR is the generation of posterior variance estimates. Future iterations should explicitly visualise these confidence intervals. This would allow a player to distinguish between a strategy that is optimal because it is ``safe'' versus one that is optimal simply because the model lacked data on the dangers of a specific area. 

Similarly, we would like to extend our framework to model par 5 holes which introduce the consideration of three chained shots in a row. Additionally, testing the model on diverse course styles ranging from links-courses (firm course with high roll and few hazards) to target golf (courses with many obstacles that force the player to fly the ball past them) would validate the robustness of our proposed framework.

Currently, the model utilizes a proxy ``average LPGA'' profile. Future work involves ingesting real shot data from distinct populations (e.g., elite male amateurs vs. seniors) to observe how optimal risk surfaces change based on power and consistency disparities.

Ultimately, this project lays the mathematical groundwork for a ``digital caddie'' evolving with the player by turning every shot recorded into a smarter decision for the future.

\appendix
\section{Data Generation Methodology}
\label{sec:Appendix}
In the absence of granular, shot-level tournament data (such as PGA Tour ShotLink data) and player-specific, we developed a robust synthetic data generation pipeline to construct realistic two-dimensional scoring distributions by combining course digitisation, empirical statistical interpolation, and probabilistic simulation.

Similarly, we constructed a proxy player profile representing an ``average'' LPGA Tour player using publicly available launch monitor data \parencite{trackman2024averages}.

\subsection{Course Geometry and Lie Classification}
\label{app: course-representation}

To create a spatially accurate environment for our model, we digitised the layout of hole 9 at Mountain Meadows Golf Course in California.
\begin{enumerate}
    \item \textbf{Vector Mapping:} The course layout was initially scraped using OpenStreetMap \parencite{openstreetmap} and imported into QGIS \parencite{qgis} for refinement.
    \item \textbf{Polygon Classification:} Features were exported as WKT (Well-Known Text) polygons, allowing for the discrete classification of any coordinate $\coord$ into specific lie types (Fairway, Rough, Bunker, Green, or Water) using the \texttt{Shapely} library in Python \parencite{shapely}. This geometric classification is the foundational layer for all subsequent scoring evaluations.
\end{enumerate}

\subsection{Green Surface Data Generation}
\label{app: green-surface data}
To model putting difficulty beyond simple distance estimates, we constructed a multivariate dataset that accounted for the complex topography of a green. This process involved a multi-stage pipeline designed to transform univariate averages into spatially aware discrete outcomes, creating data that would be more representative of that collected at an actual tournament.

\begin{enumerate}
    \item \textbf{Surface Topographical Modelling.} We defined a continuous height function $z(x_1, x_2)$ to simulate a realistic putting surface featuring a general tilt, undulating bumps, and a mid-green tier. By defining a height function, we were then able to calculate slope vectors and break severity at any point on the green.
    \item \textbf{Baseline Interpolation with Noise.} We first interpolated baseline average putts values from Mark Broadie's empirical table \ref{tab: ESHO_Putts}. 
    To mimic real-world heteroscedasticity (where variance increases with distance) and data sparsity at long range, we generated noisy synthetic samples centred at these empirical means.
\[ y_{sample} \sim \mathcal{N}(f_{Broadie}(d), \sigma_d^2) \]
We then fit an intermediate GPR model to these synthetic samples, creating a smooth, continuous interpolator capable of generating posterior predictive distributions of ESHO for any distance $d$. The Gaussian Process also captured both the mean trend and the increasing uncertainty at longer ranges.
\begin{table}[h!]
    \centering
    \small
    \caption{\textbf{PGA Tour Putting Statistics by Distance.} The ``Average putts'' column represents the PGA Tour putting benchmark used in strokes gained computations. The values are based on an analysis of nearly four million putts on the PGA Tour from 2003 to 2012, as detailed in \parencite{Broadie_2014a}.}
    \label{tab: ESHO_Putts}
    \renewcommand{\arraystretch}{1.2}
\begin{tabular}{lccccccccccccccc}
\toprule
\textbf{Distance (feet)} & 2 & 3 & 4 & 5 & 6 & 8 & 9 & 10 & 15 & 20 & 30 & 40 & 60 & 90 \\
\midrule
\textbf{One-putt \%}   & 99 & 96 & 88 & 77 & 66  & 50 & 45 & 40 & 23 & 15 & 7 & 4  & 2 & 1 \\
\textbf{Three-putt \%} & 0.0 & 0.1 & 0.3 & 0.4 & 0.4  & 0.6 & 0.7 & 0.7 & 1.3 & 2.2 & 5.0 & 10.0 & 23.0 & 41.0 \\
\textbf{Avg putts}      & 1.01 & 1.04 & 1.13 & 1.23 & 1.34 & 1.50 & 1.56 & 1.61 & 1.78 & 1.87 & 1.98 & 2.06 & 2.21 & 2.40 \\
\bottomrule
\end{tabular}

\end{table}
\item \textbf{Difficulty Multipliers}
To account for the non-linear difficulty introduced by the green's shape, we computed difficulty multipliers based on the local topography. We evaluated the gradient of the surface at four critical points along the putt's trajectory (50\%, 70\%, 85\%, and 95\% of the distance to the hole) to capture the relevant slope and break.

\begin{itemize}
    \item \textbf{Fall Line Slope:} Penalising downhill putts (where the ball picks up speed) more heavily than uphill ones.
    \item \textbf{Tier Difference:} Adding a significant penalty if the ball and hole are located on different tiers.
    \item \textbf{Side Slope:} Calculating the perpendicular slope at different points of the putt's path to account for significant break.
\end{itemize}

The specific multipliers (detailed in Table \ref{tab:multipliers}) were estimated heuristically based on domain expertise, as empirical data for slope-specific putting difficulty is not publicly available. These multipliers adjust the baseline ESHO ($\mu_{base}$) to form an adjusted mean $\mu_{adj}$:
\[ \mu_{adj} = \mu_{base} \times M_{tier} \times M_{side} \times M_{slope} \]

\begin{table}[h!]
    \centering
    \small
    \caption{\textbf{Topographical Difficulty Multipliers.} These coefficients are applied to the baseline expected strokes to hole out (ESHO) to account for green complexity. \textit{Note: These values are heuristic estimates derived from domain knowledge.}}
    \label{tab:multipliers}
    \renewcommand{\arraystretch}{1.2}
    \begin{tabular}{llc}
        \toprule
        \textbf{Factor} & \textbf{Condition} & \textbf{Multiplier Range} \\
        \midrule
        \textbf{Tier Difference} & No Tier (Diff $< 0.15$) & $1.00$ \\
        (Height diff. between ball \& hole) & Level 1 (Diff $0.15 - 0.45$) & $1.05 - 1.07$ \\
        & Level 2 (Diff $0.45 - 0.70$) & $1.08 - 1.12$ \\
        & Level 3 (Diff $0.70 - 1.00$) & $1.12 - 1.17$ \\
        & Level 4 (Diff $1.00 - 2.00$) & $1.18 - 1.25$ \\
        & Level 5 (Diff $> 2.00$) & $1.40$ \\
        \midrule
        \textbf{Side Slope} & Negligible ($< 0.5\%$) & $1.00$ \\
        (Max perpendicular slope) & Mild ($0.5 - 1.0\%$) & $1.00 - 1.03$ \\
        & Moderate ($1.0 - 2.0\%$) & $1.04 - 1.09$ \\
        & Severe ($2.0 - 3.0\%$) & $1.10 - 1.20$ \\
        & Extreme ($3.0 - 4.0\%$) & $1.20 - 1.25$ \\
        & Unplayable ($> 4.0\%$) & $1.35$ \\
        \midrule
        \textbf{Uphill / Downhill} & Flat & $1.00$ \\
        (Max parallel slope) & Uphill (Mild to Steep) & $1.03 - 1.10$ \\
        & Downhill (Mild) & $1.01 - 1.07$ \\
        & Downhill (Steep) & $1.10 - 1.45$ \\
        & Downhill (Extreme $>3\%$) & $1.45 - 1.60$ \\
        \bottomrule
    \end{tabular}
\end{table}


\item\textbf{Generating Discrete Outcomes}

To generate the final discrete stroke counts used for modeling we executed a probabilistic simulation process that transforms continuous difficulty estimates into discrete scoring distributions:
\begin{enumerate}
    \item \textbf{Sampling and Evaluation:} We selected random $\coord$ coordinates on the digitized green surface. For each point, we computed the Euclidean distance to the hole and queried the intermediate GPR (from Step 2) to obtain the posterior distribution of ESHO at that distance. We then applied the topographical multipliers (from Step 3) to this distribution, shifting the mean to reflect the specific difficulty of the putt's slope and break.

    \item \textbf{Base Probability Retrieval:} Simultaneously, we retrieved the baseline probabilities $p_1, p_2, p_{3+}$ for holing out in 1, 2, or 3+ strokes from Broadie's empirical table \ref{tab: ESHO_Putts}, based solely on the Euclidean distance to the hole.

    \item \textbf{Probability Mass Redistribution:} We utilised the shifted ESHO from the GPR to adjust these baseline probabilities. Rather than simply rounding a continuous mean, we redistributed probability mass between discrete outcomes. For example, in locations with high difficulty multipliers (e.g., steep downhill slopes or cross-tier putts), probability mass was shifted from $p_1$ to $p_2$ and $p_{3+}$, decreasing the likelihood of a one-putt while increasing the risk of a three-putt. This shift was calibrated to ensure the resulting expected value matched our adjusted ESHO ($\mu_{adj}$) while preserving the internal logic of discrete golf scoring.

    \item \textbf{Sampling and Optimization:} We sampled final discrete scores using a discrete inverse transformation on these adjusted probability vectors. Crucially, this framework allowed us to directly extract the probability of a one-putt, $P(\text{Score}=1 \mid x, y)$, from any location, and to fit the GPR in Section \ref{subsec: par4} to data that emulated the ideal data we would collect in a tournament described in \ref{ref:idealData}. 
\end{enumerate}
\end{enumerate}

\subsection{Short Game Data Generation}
\label{app: short-game data}

For shots originating off the green, we employ a direct interpolation method. Since slope is less critical for these shots than the lie condition, we estimate the expected strokes to hole out (ESHO) by interpolating directly from historical PGA Tour averages based on the distance to the hole and the specific lie type (Fairway, Rough, or Sand).

Table \ref{tab:broadie_shortgame} displays the baseline values used for this interpolation, once again derived from Broadie's work \parencite{Broadie_2014a}

\begin{table}[h!]
    \centering
    \small
    \caption{\textbf{Average Shots to Hole Out by Lie and Distance.} Values represent the expected number of strokes to hole out from specific distances and lies, derived from PGA Tour ShotLink data (2003-2012) \parencite{Broadie_2014a}.}
    \label{tab:broadie_shortgame}
    \renewcommand{\arraystretch}{1.2}
    \begin{tabular}{cccc}
        \toprule
        \textbf{Distance} & \textbf{Fairway} & \textbf{Rough} & \textbf{Sand} \\
        \textbf{(yards)} & \textbf{(avg shots)} & \textbf{(avg shots)} & \textbf{(avg shots)} \\
        \midrule
        20 & 2.40 & 2.59 & 2.53 \\
        40 & 2.60 & 2.78 & 2.82 \\
        60 & 2.70 & 2.91 & 3.15 \\
        80 & 2.75 & 2.96 & 3.24 \\
        \bottomrule
    \end{tabular}
\end{table}

\subsubsection{Simulating Player-Specific Data}
\label{app: player-data}
We constructed a proxy player profile representing an ``average'' right-handed LPGA Tour player using a simulation based on publicly available carry distance averages \parencite{trackman2024averages}.
\begin{enumerate}\item 
\textbf{Base Averages:} We defined average carry distances for 20 distinct shot types, ranging from a Driver ($223$ yards) to a $60^{\circ}$ wedge ($78$ yards) based off of official data; also including partial wedge shots to bridge yardage gaps for the shorter distances (e.g., a $3/4$ 60-degree wedge). 
\item \textbf{Dispersion Scaling:} To model the natural increase in variance with shot length, we scaled the standard deviations for lateral error ($\sigma_{side}$) and distance error ($\sigma_{dist}$) linearly with carry distance, setting $\sigma_{side} \approx 0.08 \times \text{Avg. Carry}$ and $\sigma_{dist} \approx 0.04 \times \text{Avg. Carry}$. 
\item \textbf{Directional Bias:} We introduced a correlation between lateral miss and carry distance to mimic real ball flight physics. Shots missing left (simulating pulls) were modelled with a positive distance bias ($+5$ yards on average), while shots missing right (simulating fades or pushes) were modelled with a negative distance bias ($-2$ yards on average), this adjustment was heuristically based on domain expertise, as empirical data for these distributions is not publicly available.
\item \textbf{Sampling:} For each club, we generated $N=30$ synthetic shot outcomes by sampling from these distributions, creating the dataset used to fit the player's bivariate Gaussian profiles.
\end{enumerate}

\section{Glossary of Terms}
\label{sec:Appendix2}

\begin{itemize}
    \item \textbf{Lie:} The ground condition or surface where a ball rests. Common lies include:
    \begin{itemize}
        \item\textit{Fairway}: Closely mown grass offering the standard hitting surface.
        \item\textit{Rough}: Longer grass bordering fairway lies, thicker and more unpredictable.
        \item\textit{Sand/bunker}: Obstacle or hazard filled with sand.
        \item\textit{Green}: Short grass area around the hole used for putting.
    \end{itemize}
    \item\textbf{Club:} Instrument used to hit golf balls. Clubs vary in loft and length, affecting distance and shot shape. Examples used in this paper include: 5-iron, Wood or Driver.
    \item\textbf{Tee/Tee box:} Designated starting area for each hole.
    \item\textbf{Pin/Flag:} The flagstick marking the hole's location on the green.
    \item\textbf{Hazard:} Course obstacles like water or sand bunkers that penalise errant shots.
    \item\textbf{Drop:} Placing the ball back in play after it lands in a water hazard, typically incurring a penalty stroke.
    \item \textbf{Par:} The expected number of strokes a skilled golfer should take to complete a hole. A par-3 hole expects a golfer to finish in 3 shots.
    \item \textbf{Birdie:} The score that is one stroke under par on a hole.
    \item \textbf{Bogey:} The score that is one stroke over par on a hole.
    \item\textbf{Stroke:} A single swing attempt at the ball.
    \item\textbf{Stock shot:} A player's standard, repeatable ``full swing'' given a club.
    \item\textbf{Approach shot:} A shot intended to reach the green, or ``approach" the pin.
    \item \textbf{Carry distance:} The distance a ball travels through the air before it first touches the ground.
    \item \textbf{Dispersion:} The spread of a player’s shots around their intended target, variable depending on golf club difficulty and individual golfer's tendencies.
    \item \textbf{Aimpoint:} Directional target chosen by the player, defining the centre of their shot dispersion.
    \item \textbf{Short game:} Shots played close to the green (chipping, pitching, putting).
    \item \textbf{Break:} The curve a putt takes due to green slope.
    \item \textbf{Trackman:} A radar launch monitor system that captures detailed shot data, including launch angle, ball speed, spin, and dispersion.
    \item \textbf{ShotLink:} The PGA Tour's proprietary real-time scoring system that captures 3-dimensional coordinates for every shot taken during a tournament.
    \item \textbf{Strokes Gained (SG):} A statistical metric that compares a player's performance against a baseline (typically the PGA Tour average) to quantify the quality of a shot.
    \item \textbf{Smash Factor:} The ratio of ball speed to clubhead speed, used as a measure of energy transfer efficiency.
    \item \textbf{WAGR:} World Amateur Golf Ranking.
\end{itemize}











% -------------------------------------------------------
% END MATTER
% -------------------------------------------------------
\newpage 



\printbibliography[title={References}]

\vspace{2cm}


\section*{Generative AI Statement}
\noindent
Brief statement describing how (if at all) generative AI tools were used in your work and the preparation of your submission.

\end{document}
